\documentclass[12pt]{article}
\usepackage[utf8]{inputenc}
\usepackage[spanish]{babel}
\usepackage{amsmath}
\usepackage{amssymb}
\usepackage{geometry}
\geometry{a4paper, margin=1in}
\begin{document}
\section{Arquitectura del Código y Flujo de Trabajo}

La arquitectura del simulador físico Cosmic-Generator está estructurada para maximizar la modularidad, el rendimiento y la facilidad de uso. El flujo de trabajo o pipeline sigue un modelo secuencial desde la captura de parámetros del usuario hasta la visualización en tiempo real y la síntesis de audio.

\subsection{Punto de Entrada: UI y Configuración (\texttt{app.py})}
El ciclo de ejecución de la plataforma comienza en el archivo \texttt{app.py}. Este script sirve como punto de entrada de la aplicación y utiliza el framework Streamlit para proveer una Interfaz de Usuario (UI) web interactiva. En esta etapa, \texttt{app.py} permite al usuario definir las variables globales y parámetros específicos del sistema (tales como constantes físicas, niveles de fricción, propiedades visuales, entre otros). Al iniciar la simulación, los parámetros se empaquetan y se transmiten al siguiente nivel de la arquitectura.

\subsection{Orquestación del Sistema (\texttt{lanzador.py})}
El módulo \texttt{lanzador.py} cumple el rol de orquestador o controlador principal. Recibe los parámetros configurados en la UI e inicializa todas las piezas fundamentales del simulador. Sus responsabilidades incluyen:
\begin{itemize}
    \item Instanciar el entorno físico.
    \item Configurar los subsistemas de generación de audio y renderizado visual.
    \item Disparar el bucle principal de simulación, delegando la carga computacional a los módulos especializados.
\end{itemize}
De esta forma, se mantiene completamente desacoplada la interfaz web de la lógica pesada del simulador.

\subsection{Patrón Productor-Consumidor en Renderizado (\texttt{src/render/stable/render\_standard.py})}
El rendimiento es crítico en simulaciones físicas de tiempo real. Para asegurar una alta tasa de fotogramas por segundo (FPS), el archivo \texttt{src/render/stable/render\_standard.py} implementa un modelo de concurrencia basado en el patrón \textbf{Productor-Consumidor}, empleando múltiples hilos (threads):

\begin{itemize}
    \item \textbf{Productor (Hilo Físico):} Un hilo en segundo plano se dedica exclusivamente a resolver las ecuaciones matemáticas del simulador. Calcula la cinemática, actualiza posiciones, velocidades y detecta colisiones, produciendo \textit{estados del sistema} que se almacenan de forma segura en una estructura de datos concurrente (cola o buffer).
    \item \textbf{Consumidor (Hilo Gráfico):} El hilo principal de renderizado lee o "consume" los estados físicos ya calculados que están disponibles en el buffer, para dibujar cada fotograma. Al separar estos dos procesos, el cálculo matemático no bloquea el repintado de la pantalla y la visualización gráfica no ralentiza la física, logrando una aceleración notable de los FPS.
\end{itemize}

\subsection{Organización Modular del Código}
El código fuente se organiza de forma rigurosa y cohesiva en distintos directorios bajo \texttt{src/}, respetando el principio de responsabilidad única:

\begin{itemize}
    \item \textbf{\texttt{src/core/}:} Contiene el motor numérico y físico del proyecto. Aquí residen los cálculos de dinámica gravitacional, cinemática y detección de colisiones. Este módulo opera con matemática pura y no tiene conocimiento del entorno gráfico ni sonoro.
    \item \textbf{\texttt{src/audio/}:} Módulo encargado de la síntesis y gestión de sonido. Actúa escuchando los eventos disparados por el Core (por ejemplo, el rebote o impacto de una entidad) y los transforma dinámicamente en efectos de audio espaciales o arreglos musicales.
    \item \textbf{\texttt{src/render/}:} Dedicado de forma exclusiva a la representación visual. Recopila la información geométrica del Core y realiza las llamadas a la biblioteca gráfica pertinente para renderizar el entorno y los elementos simulados en la ventana de la aplicación.
\end{itemize}

\end{document}
